\chapter{操作系统绪论}

\section{操作系统的概念与定位}
\concept{计算机系统的层次}
\begin{points}
  \item 计算机系统由\term{硬件}和\term{软件}组成。硬件包括 CPU、内存、I/O 设备等；软件包括系统软件和应用软件。
  \item 操作系统是配置在裸机上的\key{第一层软件}，位于硬件与其他软件之间，是其他软件运行的基础。
  \item 编译器、数据库管理系统、编辑程序也属于系统软件，但它们不是操作系统本身。
  \item 从层次上看：硬件/裸机 $\rightarrow$ 操作系统 $\rightarrow$ 系统程序 $\rightarrow$ 应用程序 $\rightarrow$ 用户。
\end{points}
\plain{硬件不会直接服务应用，OS 是硬件上第一层“扩展机”，把裸机包装成可管理、可调用的机器。}
\exam{常考“谁是第一层软件”“编译器/数据库是不是 OS”“OS 位于哪两层之间”。}

\concept{为什么需要操作系统}
\begin{points}
  \item \term{从用户视角看}：操作系统是用户和计算机硬件之间的接口，为用户提供方便、安全、高效的使用环境。
  \item \term{从系统视角看}：操作系统是计算机系统资源的管理者，负责分配、回收、保护和调度资源。
  \item 操作系统通过三种思想扩展计算机能力：\key{资源复用、资源虚化、资源抽象}。
  \item 复用解决“资源数量不够”：CPU 轮流给多个进程用是时间复用；内存分块给多个进程用是空间复用。
  \item 虚化解决“一个物理资源变成多个逻辑资源”：一个 CPU 虚拟成多个虚拟处理机，一台打印机通过 SPOOLing 虚拟成多台逻辑打印机。
  \item 抽象解决“硬件太复杂”：磁盘扇区被抽象成文件，CPU 执行流被抽象成进程，物理设备被抽象成逻辑设备名。
\end{points}
\plain{没有 OS，用户要直接面对 CPU、内存、磁盘、设备；OS 的价值就是管资源、挡危险、把复杂硬件变简单。}
\exam{常考从用户视角/系统视角回答，或解释资源复用、虚化、抽象。}


\concept{操作系统定义}
\begin{points}
  \item 操作系统没有唯一的精确定义，但可从两个角度理解。
  \item \term{控制程序}：控制程序执行，防止错误和不当使用，协调应用与硬件。
  \item \term{资源管理器}：管理处理机、存储器、I/O 设备、文件、磁盘空间等软硬件资源。
  \item 简洁表述：操作系统是位于裸机和应用程序之间的系统软件，对下管理控制计算机软硬件资源，对上提供方便、安全、高效的接口和运行环境。
\end{points}
\plain{定义题不要追求一句神句，抓住“系统软件、资源管理者、用户接口/扩展机”三个关键词就稳。}
\exam{常考简答“什么是操作系统”，按“位置-对象-作用-接口”写。}


\section{操作系统的形成与发展}
\concept{发展阶段总览}
\begin{points}
  \item 操作系统发展顺序：\key{手工操作 $\rightarrow$ 批处理系统 $\rightarrow$ 多道批处理系统 $\rightarrow$ 分时系统 $\rightarrow$ 实时系统 $\rightarrow$ 通用操作系统}。
  \item 每个阶段都是在解决上一阶段的主要矛盾：人机矛盾、CPU 与 I/O 速度矛盾、交互需求、实时响应需求。
\end{points}
\plain{发展史就是不断解决矛盾：人机慢、CPU 等 I/O、用户要交互、任务有截止期限。}
\exam{常考按阶段排序，或问某阶段解决了什么问题。}


\concept{手工操作阶段}
\begin{points}
  \item 用户独占机器，人工装入程序、启动运行、取走结果。
  \item 主要矛盾：人太慢，CPU 太快，I/O 太慢，导致资源利用率很低。
  \item 没有真正的操作系统，只是人手工控制机器。
\end{points}
\plain{人手工装程序，CPU 大量等人和 I/O，所以效率极低。}
\exam{常考缺点：用户独占资源、CPU 利用率低、没有自动作业控制。}


\concept{批处理系统}
\begin{points}
  \item 把多个作业组织成一批，由监督程序 Monitor 自动控制，一个接一个执行。
  \item 优点：减少人工干预，提高吞吐量。
  \item 缺点：缺少交互性，用户提交作业后不能直接干预运行过程。
  \item 关键词：\key{自动、成批、无交互、吞吐量、周转时间}。
\end{points}
\plain{把一堆作业交给监督程序自动跑，减少人工干预，但用户不能边运行边改。}
\exam{常考关键词“成批、自动、无交互、吞吐量”。}


\concept{多道批处理系统}
\begin{points}
  \item 多个作业同时进入内存，交替运行；当一个作业等待 I/O 时，CPU 可执行另一个作业。
  \item 核心目标：解决 CPU 等待 I/O 的浪费，提高 CPU 和 I/O 设备利用率。
  \item 代价：需要进程调度、内存分配、资源保护、中断处理、上下文切换等机制，系统开销更大。
  \item 直觉例子：作业 A 等磁盘读入时，作业 B 用 CPU 计算；磁盘和 CPU 可以并行工作。
\end{points}
\plain{多个作业进内存，一个等 I/O，另一个用 CPU，让 CPU 和设备都别闲着。}
\exam{常考为什么能提高效率、为什么需要中断/调度/存储保护。}


\concept{分时系统}
\begin{points}
  \item 多个用户通过终端同时使用计算机，系统把 CPU 时间划分为时间片，轮流响应各用户请求。
  \item 追求目标：较快响应用户，使每个用户感觉自己独占机器。
  \item 特征：\key{多路性、交互性、独立性、及时性}。
  \item 注意：分时系统的及时性是“用户命令较快响应”，不是实时系统那种严格截止期限。
  \item 若时间片为 $Q$，用户数为 $N$，粗略响应周期可理解为 $T=Q\times N$。用户越多，每个用户轮到 CPU 的间隔越长。
\end{points}
\plain{把 CPU 切成很多时间片轮流给用户，目标是让每个终端都感觉机器在响应自己。}
\exam{常考分时四特征，尤其“及时性”不等于硬实时。}


\concept{实时系统}
\begin{points}
  \item 实时系统强调在规定时间内响应外部事件。不是单纯“速度快”，而是响应必须\key{及时、可预测、可靠}。
  \item 硬实时：错过截止时间会造成严重失败，例如飞控、安全控制。
  \item 软实时：偶尔超时可以容忍，例如多媒体播放中偶尔卡顿。
  \item 关键词：\key{及时性、可靠性、截止时间、硬实时、软实时}。
\end{points}
\plain{实时不是“跑得快”而是“截止时间前必须响应”。}
\exam{常考硬实时/软实时区别、实时系统追求可靠性和可预测性。}


\section{操作系统的基本特征}
\concept{并发性}
\begin{points}
  \item 并发指若干事件在\key{同一时间间隔内}发生。
  \item 并行指若干事件在\key{同一时刻}发生。
  \item 单 CPU 上多个进程可以并发，但不能并行；多 CPU/多核上多个进程可以真正并行。
  \item 课件常用判断：并行一定是并发，但并发不一定是并行；并行是并发的特例。
\end{points}
\plain{核心是“看起来同时推进”。单 CPU 上靠时间片交替，多 CPU 上才可能真的同一时刻运行。}
\exam{常考并发与并行区别、单处理机下最多几个进程运行、并发导致为什么需要同步互斥。}


\concept{共享性}
\begin{points}
  \item 共享指系统资源可供多个并发执行的进程共同使用。
  \item \term{互斥共享}：一段时间内只允许一个进程访问，例如打印机、临界区。
  \item \term{同时共享}：宏观上多个进程可同时访问，例如只读文件、共享代码段。
  \item 并发与共享互为条件：没有并发，就谈不上资源竞争；没有共享管理，并发程序也难以正确运行。
\end{points}
\plain{多个进程都想用同一批资源，所以必须规定能不能一起用、谁先用、用完怎么释放。}
\exam{常考互斥共享/同时共享举例，或者问并发和共享为什么是 OS 最基本特征。}


\concept{虚拟性}
\begin{points}
  \item 虚拟性指通过某种技术把一个物理实体变成若干逻辑上的对应物。
  \item 典型例子：虚拟处理机、虚拟存储器、虚拟设备。
  \item 虚拟不是“凭空变多”，而是通过复用、映射、换入换出等机制，让用户看到更方便的逻辑资源。
\end{points}
\plain{不是资源真的变多，而是 OS 用映射、换入换出、排队等办法，让用户感觉自己有一份独立资源。}
\exam{常考虚拟处理机、虚拟存储器、虚拟设备的例子，以及虚拟性和资源复用的关系。}


\concept{异步性/不确定性}
\begin{points}
  \item 异步性指并发环境下，程序以不可预知的速度向前推进。
  \item 一个进程何时被调度、何时暂停、I/O 等多久、何时继续执行，不完全由程序自己决定。
  \item 异步性带来问题：若没有同步互斥机制，可能出现数据不一致、结果不可再现。
\end{points}
\plain{进程什么时候跑、什么时候停，不完全由程序自己决定，而是被调度、中断、I/O 共同影响。}
\exam{常考并发程序为什么不可再现、为什么会有与时间有关的错误。}


\section{操作系统的功能与接口}
\concept{功能模块}
\begin{points}
  \item \term{处理机管理/进程管理}：进程控制、进程同步、进程通信、调度。
  \item \term{存储器管理}：内存分配与回收、地址映射、保护共享、存储扩充。
  \item \term{设备管理}：I/O 控制、设备分配、缓冲、设备驱动、设备独立性。
  \item \term{文件管理}：文件存储、目录、文件共享、文件保护、外存空间管理。
  \item \term{作业管理/用户接口}：作业控制、命令接口、系统调用接口、图形接口。
\end{points}
\plain{五大管理就是 OS 的主线：CPU、内存、设备、文件、用户/作业。后面章节基本都从这里展开。}
\exam{常考某功能属于哪个模块，如进程调度属于处理机管理，文件目录属于文件管理。}


\concept{操作系统提供的接口}
\begin{points}
  \item \term{命令接口}：用户直接输入命令控制系统，如 \code{ls}、\code{cd}、\code{mkdir}。Shell 本质是命令解释器。
  \item 命令接口分为联机用户接口和脱机用户接口。联机接口用于交互式控制；脱机接口常见于批处理作业控制语言。
  \item \term{程序接口}：程序通过系统调用请求 OS 服务，如 \code{open()}、\code{read()}、\code{fork()}。
  \item \term{图形接口}：窗口、菜单、鼠标等 GUI 形式，本质上也是调用系统服务。
  \item 广义指令常作为系统调用的另一种说法，属于程序接口。
\end{points}
\plain{人用命令/GUI 找 OS，程序用系统调用找 OS；危险操作必须让内核代办。}
\exam{常考命令接口、程序接口、GUI 的分类，以及系统调用是否属于程序接口。}

\concept{系统调用与特权操作}
\begin{points}
  \item I/O 指令、修改中断开关、修改页表、设置定时器等通常是特权指令，只能在内核态执行。
  \item 普通用户程序不能直接访问设备或随意操作内存，否则可能破坏其他进程或系统本身。
  \item 用户程序通过系统调用陷入内核，切换到内核态，由操作系统执行受保护的操作，再返回用户态。
\end{points}
\plain{人用命令/GUI 找 OS，程序用系统调用找 OS；危险操作必须让内核代办。}
\exam{常考命令接口、程序接口、GUI 的分类，以及系统调用是否属于程序接口。}


\section{操作系统结构}
\concept{常见结构}
\begin{points}
  \item \term{简单结构}：系统各部分边界不清，效率可能高，但维护困难。
  \item \term{层次结构}：从底层硬件到上层服务逐层构造，结构清晰，但层次划分和跨层调用可能带来开销。
  \item \term{微内核结构}：内核只保留最基本功能，如进程通信、低级地址空间管理、基本调度；文件系统、设备驱动等尽量放到用户态服务中。
  \item \term{模块化结构}：核心内核加可加载模块，兼顾性能与可扩展性。
  \item \term{虚拟机结构}：虚拟机监控器在一台物理机上抽象出多台虚拟计算机，每台虚拟机可运行自己的 OS。
\end{points}
\plain{OS 结构是在“好维护”和“高效率”之间权衡：简单结构快但乱，层次清楚但可能慢，微内核可靠但通信开销大。}
\exam{常考宏内核/微内核/模块化优缺点。}


